% !TEX root = ../main.tex

\chapter{总结与展望}

\section{主要完成工作}

本文围绕“上游速度脉动对双喷嘴横向射流火焰熄火边界的影响”这一核心问题，从实验平台搭建、声激励系统研制、双色PIV诊断方法建立到前置实验开展，完成了以下主要工作：

（1）\textbf{横向射流火焰实验平台搭建}：完成了两级串联横向射流燃烧实验台的搭建与调试。一级燃烧室可产生温度范围为1200--1520~K的高温产物气流，经收缩段均匀化后进入二级燃烧室；二级燃烧室采用石英玻璃材质，便于光学诊断与高速摄影观测。实验台整体采用铝型材框架支撑，具备良好的移动性与光学对准便利性。

（2）\textbf{双喷嘴设计与加工}：设计并加工了四组不同孔径（2~mm/3~mm）与间距（2$d$/3$d$）的双喷嘴，明确了316L不锈钢材质、内壁光滑度Ra≤1.6等加工要求。6~mm孔径双喷嘴已用于预实验与系统调试，2~mm与3~mm孔径双喷嘴已完成加工，将用于正式实验。

（3）\textbf{声激励脉动系统研制}：成功研制出两套功能完整、气密性良好的声激励罐体，分别用于甲烷与空气气路的独立激励。系统可实现45~Hz--2.5~kHz频率范围内正弦波与方波激励的稳定输出，扬声器安装于罐体尾部、气体从侧壁引入的设计有效避免了气流对扬声器膜片的冲击。经密封测试与电路布局优化，系统运行稳定可靠。

（4）\textbf{双色PIV诊断系统建立}：完成了双色PIV光路系统的搭建，采用365~nm紫外LED阵列作为光源，配合正交柱面镜组与光学狭缝，形成了厚度约1--2~mm的片光。建立了从原始双色图像到速度场与混合物分数的全流程数据处理方法，实现了红蓝粒子的自适应分离、速度场求解与混合物分数计算。核心技术已申请发明专利（申请号：LSBJ-25-I-91584335）。此外，采用DMD方法对冷态流场进行了模态分析，识别了流场的主导空间结构与时间节律，为后续声激励工况的对比分析提供了基准。

（5）\textbf{前置实验与工况筛选}：完成了一级燃烧室多工况火焰燃烧实验，筛选出5组现阶段实验工况，出口温度覆盖1200--1500~K，为二级实验提供了可靠的热来流条件。完成了二级燃烧室横向射流火焰的预实验，采用乙烯/空气进行现象观测，验证了实验系统的基本功能。预实验发现6~mm孔径喷嘴存在燃烧不充分、观察窗熏黑的问题，明确了后续需采用2~mm/3~mm小孔径双喷嘴的改进方向。

（6）\textbf{激励传递特性初步评估}：采用热膜风速仪对声激励脉动系统的激励效果进行了初步测试。结果表明，热膜风速仪无法有效响应高频激励信号，未能测得有效的速度脉动幅值。后续将采用高频PIV或双色PIV速度场数据直接计算脉动幅值，以获取声激励传递特性的定量数据。

\section{创新点}

本研究将采用双色粒子图像测速（PIV）等先进燃烧诊断技术，建立流场与火焰特性的同步测量系统。通过该系统可同步获取横向射流场的速度分布、湍流强度等流场信息，以及火焰温度分布、火焰面形态、燃烧反应区演化等火焰特性参数，基于多维度同步数据的耦合分析，实现对声激励调控火焰熄火过程机理层面的深度解析与本质洞察。

本研究聚焦于单喷嘴及双喷嘴构型下横向射流火焰的熄火特性，系统探究声激励参数对火焰熄火过程的调控规律，并剖析双喷嘴布局中喷嘴间流场干涉、燃料混合及火焰相互作用的内在机理。

研究核心目标在于突破传统燃烧特性研究中以现象观测为核心的局限，通过主动调控声激励的频率、振幅等关键参数，通过调节声激励脉动系统的各项参数，进而筛选并确定能够产生正效应（即显著扩展火焰熄火极限、提升燃烧稳定性）的优化声学参数区间，为燃烧系统的主动控制技术研发提供理论支撑与参数依据。


\section{下一步工作计划}

（1）\textbf{小孔径双喷嘴正式实验}：采用2~mm与3~mm孔径双喷嘴，系统开展无声激励条件下的横向射流火焰熄火边界测量。重点探究孔径与间距对熄火当量比、临界射流速度的影响规律，确定最优喷嘴构型。

（2）\textbf{声激励下熄火边界测量}：在最优喷嘴构型基础上，施加不同频率（50--1000~Hz）、波形（正弦波/方波）与幅值的声激励，系统测量有声激励条件下的熄火边界，揭示声激励参数对熄火极限的调控规律。

（3）\textbf{双色PIV流场-混合协同诊断}：在典型工况下（无激励、最优激励、过强激励），开展双色PIV测量，同步获取燃料流与空气流的速度场及混合物分数分布，从流场-混合耦合视角揭示声激励致熄的物理机理。

（4）\textbf{脉动幅值精确测量}：针对热膜风速仪未能有效测量的问题，采用高频PIV或基于双色PIV速度场数据的后处理方法，获取声激励下喷嘴出口附近的速度脉动幅值与频谱特性，建立“激励信号→速度脉动”的定量关联。

（5）\textbf{数据整理与论文撰写}：整合全部实验数据，完成速度脉动-混合物分数-火焰熄火的量化关联模型，撰写毕业论文初稿。